% Options for packages loaded elsewhere
\PassOptionsToPackage{unicode}{hyperref}
\PassOptionsToPackage{hyphens}{url}
\PassOptionsToPackage{dvipsnames,svgnames,x11names}{xcolor}
%
\documentclass[
  letterpaper,
  DIV=11,
  numbers=noendperiod]{scrreprt}

\usepackage{amsmath,amssymb}
\usepackage{iftex}
\ifPDFTeX
  \usepackage[T1]{fontenc}
  \usepackage[utf8]{inputenc}
  \usepackage{textcomp} % provide euro and other symbols
\else % if luatex or xetex
  \usepackage{unicode-math}
  \defaultfontfeatures{Scale=MatchLowercase}
  \defaultfontfeatures[\rmfamily]{Ligatures=TeX,Scale=1}
\fi
\usepackage{lmodern}
\ifPDFTeX\else  
    % xetex/luatex font selection
\fi
% Use upquote if available, for straight quotes in verbatim environments
\IfFileExists{upquote.sty}{\usepackage{upquote}}{}
\IfFileExists{microtype.sty}{% use microtype if available
  \usepackage[]{microtype}
  \UseMicrotypeSet[protrusion]{basicmath} % disable protrusion for tt fonts
}{}
\makeatletter
\@ifundefined{KOMAClassName}{% if non-KOMA class
  \IfFileExists{parskip.sty}{%
    \usepackage{parskip}
  }{% else
    \setlength{\parindent}{0pt}
    \setlength{\parskip}{6pt plus 2pt minus 1pt}}
}{% if KOMA class
  \KOMAoptions{parskip=half}}
\makeatother
\usepackage{xcolor}
\setlength{\emergencystretch}{3em} % prevent overfull lines
\setcounter{secnumdepth}{-\maxdimen} % remove section numbering
% Make \paragraph and \subparagraph free-standing
\makeatletter
\ifx\paragraph\undefined\else
  \let\oldparagraph\paragraph
  \renewcommand{\paragraph}{
    \@ifstar
      \xxxParagraphStar
      \xxxParagraphNoStar
  }
  \newcommand{\xxxParagraphStar}[1]{\oldparagraph*{#1}\mbox{}}
  \newcommand{\xxxParagraphNoStar}[1]{\oldparagraph{#1}\mbox{}}
\fi
\ifx\subparagraph\undefined\else
  \let\oldsubparagraph\subparagraph
  \renewcommand{\subparagraph}{
    \@ifstar
      \xxxSubParagraphStar
      \xxxSubParagraphNoStar
  }
  \newcommand{\xxxSubParagraphStar}[1]{\oldsubparagraph*{#1}\mbox{}}
  \newcommand{\xxxSubParagraphNoStar}[1]{\oldsubparagraph{#1}\mbox{}}
\fi
\makeatother

\usepackage{color}
\usepackage{fancyvrb}
\newcommand{\VerbBar}{|}
\newcommand{\VERB}{\Verb[commandchars=\\\{\}]}
\DefineVerbatimEnvironment{Highlighting}{Verbatim}{commandchars=\\\{\}}
% Add ',fontsize=\small' for more characters per line
\usepackage{framed}
\definecolor{shadecolor}{RGB}{241,243,245}
\newenvironment{Shaded}{\begin{snugshade}}{\end{snugshade}}
\newcommand{\AlertTok}[1]{\textcolor[rgb]{0.68,0.00,0.00}{#1}}
\newcommand{\AnnotationTok}[1]{\textcolor[rgb]{0.37,0.37,0.37}{#1}}
\newcommand{\AttributeTok}[1]{\textcolor[rgb]{0.40,0.45,0.13}{#1}}
\newcommand{\BaseNTok}[1]{\textcolor[rgb]{0.68,0.00,0.00}{#1}}
\newcommand{\BuiltInTok}[1]{\textcolor[rgb]{0.00,0.23,0.31}{#1}}
\newcommand{\CharTok}[1]{\textcolor[rgb]{0.13,0.47,0.30}{#1}}
\newcommand{\CommentTok}[1]{\textcolor[rgb]{0.37,0.37,0.37}{#1}}
\newcommand{\CommentVarTok}[1]{\textcolor[rgb]{0.37,0.37,0.37}{\textit{#1}}}
\newcommand{\ConstantTok}[1]{\textcolor[rgb]{0.56,0.35,0.01}{#1}}
\newcommand{\ControlFlowTok}[1]{\textcolor[rgb]{0.00,0.23,0.31}{\textbf{#1}}}
\newcommand{\DataTypeTok}[1]{\textcolor[rgb]{0.68,0.00,0.00}{#1}}
\newcommand{\DecValTok}[1]{\textcolor[rgb]{0.68,0.00,0.00}{#1}}
\newcommand{\DocumentationTok}[1]{\textcolor[rgb]{0.37,0.37,0.37}{\textit{#1}}}
\newcommand{\ErrorTok}[1]{\textcolor[rgb]{0.68,0.00,0.00}{#1}}
\newcommand{\ExtensionTok}[1]{\textcolor[rgb]{0.00,0.23,0.31}{#1}}
\newcommand{\FloatTok}[1]{\textcolor[rgb]{0.68,0.00,0.00}{#1}}
\newcommand{\FunctionTok}[1]{\textcolor[rgb]{0.28,0.35,0.67}{#1}}
\newcommand{\ImportTok}[1]{\textcolor[rgb]{0.00,0.46,0.62}{#1}}
\newcommand{\InformationTok}[1]{\textcolor[rgb]{0.37,0.37,0.37}{#1}}
\newcommand{\KeywordTok}[1]{\textcolor[rgb]{0.00,0.23,0.31}{\textbf{#1}}}
\newcommand{\NormalTok}[1]{\textcolor[rgb]{0.00,0.23,0.31}{#1}}
\newcommand{\OperatorTok}[1]{\textcolor[rgb]{0.37,0.37,0.37}{#1}}
\newcommand{\OtherTok}[1]{\textcolor[rgb]{0.00,0.23,0.31}{#1}}
\newcommand{\PreprocessorTok}[1]{\textcolor[rgb]{0.68,0.00,0.00}{#1}}
\newcommand{\RegionMarkerTok}[1]{\textcolor[rgb]{0.00,0.23,0.31}{#1}}
\newcommand{\SpecialCharTok}[1]{\textcolor[rgb]{0.37,0.37,0.37}{#1}}
\newcommand{\SpecialStringTok}[1]{\textcolor[rgb]{0.13,0.47,0.30}{#1}}
\newcommand{\StringTok}[1]{\textcolor[rgb]{0.13,0.47,0.30}{#1}}
\newcommand{\VariableTok}[1]{\textcolor[rgb]{0.07,0.07,0.07}{#1}}
\newcommand{\VerbatimStringTok}[1]{\textcolor[rgb]{0.13,0.47,0.30}{#1}}
\newcommand{\WarningTok}[1]{\textcolor[rgb]{0.37,0.37,0.37}{\textit{#1}}}

\providecommand{\tightlist}{%
  \setlength{\itemsep}{0pt}\setlength{\parskip}{0pt}}\usepackage{longtable,booktabs,array}
\usepackage{calc} % for calculating minipage widths
% Correct order of tables after \paragraph or \subparagraph
\usepackage{etoolbox}
\makeatletter
\patchcmd\longtable{\par}{\if@noskipsec\mbox{}\fi\par}{}{}
\makeatother
% Allow footnotes in longtable head/foot
\IfFileExists{footnotehyper.sty}{\usepackage{footnotehyper}}{\usepackage{footnote}}
\makesavenoteenv{longtable}
\usepackage{graphicx}
\makeatletter
\newsavebox\pandoc@box
\newcommand*\pandocbounded[1]{% scales image to fit in text height/width
  \sbox\pandoc@box{#1}%
  \Gscale@div\@tempa{\textheight}{\dimexpr\ht\pandoc@box+\dp\pandoc@box\relax}%
  \Gscale@div\@tempb{\linewidth}{\wd\pandoc@box}%
  \ifdim\@tempb\p@<\@tempa\p@\let\@tempa\@tempb\fi% select the smaller of both
  \ifdim\@tempa\p@<\p@\scalebox{\@tempa}{\usebox\pandoc@box}%
  \else\usebox{\pandoc@box}%
  \fi%
}
% Set default figure placement to htbp
\def\fps@figure{htbp}
\makeatother

\KOMAoption{captions}{tableheading}
\makeatletter
\@ifpackageloaded{caption}{}{\usepackage{caption}}
\AtBeginDocument{%
\ifdefined\contentsname
  \renewcommand*\contentsname{Table of contents}
\else
  \newcommand\contentsname{Table of contents}
\fi
\ifdefined\listfigurename
  \renewcommand*\listfigurename{List of Figures}
\else
  \newcommand\listfigurename{List of Figures}
\fi
\ifdefined\listtablename
  \renewcommand*\listtablename{List of Tables}
\else
  \newcommand\listtablename{List of Tables}
\fi
\ifdefined\figurename
  \renewcommand*\figurename{Figure}
\else
  \newcommand\figurename{Figure}
\fi
\ifdefined\tablename
  \renewcommand*\tablename{Table}
\else
  \newcommand\tablename{Table}
\fi
}
\@ifpackageloaded{float}{}{\usepackage{float}}
\floatstyle{ruled}
\@ifundefined{c@chapter}{\newfloat{codelisting}{h}{lop}}{\newfloat{codelisting}{h}{lop}[chapter]}
\floatname{codelisting}{Listing}
\newcommand*\listoflistings{\listof{codelisting}{List of Listings}}
\makeatother
\makeatletter
\makeatother
\makeatletter
\@ifpackageloaded{caption}{}{\usepackage{caption}}
\@ifpackageloaded{subcaption}{}{\usepackage{subcaption}}
\makeatother

\usepackage{bookmark}

\IfFileExists{xurl.sty}{\usepackage{xurl}}{} % add URL line breaks if available
\urlstyle{same} % disable monospaced font for URLs
\hypersetup{
  colorlinks=true,
  linkcolor={blue},
  filecolor={Maroon},
  citecolor={Blue},
  urlcolor={Blue},
  pdfcreator={LaTeX via pandoc}}


\author{}
\date{}

\begin{document}


\chapter{Exercise solutions}\label{exercise-solutions}

\subsection*{Week 1}\label{week-1}
\addcontentsline{toc}{subsection}{Week 1}

The following exercise will allow you to test yourself against what you
have learned so far. The solutions will be released at the end of the
week.

Using the dataset
\href{https://raw.githubusercontent.com/atpinto/RM1/refs/heads/main/Data/hers_subset.csv}{hers\_subset.csv}
dataset, use simple linear regression in R or Stata to measure the
association between diastolic blood pressure (DBP - the outcome) and
body mass index (BMI - the exposure).

\subsubsection{\texorpdfstring{a) Summarise the important findings by
interpreting the relevant parameter values, associated P-values and
confidence intervals, and \(R^2\) value. Three to four sentences is
usually enough
here.}{a) Summarise the important findings by interpreting the relevant parameter values, associated P-values and confidence intervals, and R\^{}2 value. Three to four sentences is usually enough here.}}\label{a-summarise-the-important-findings-by-interpreting-the-relevant-parameter-values-associated-p-values-and-confidence-intervals-and-r2-value.-three-to-four-sentences-is-usually-enough-here.}

\subsubsection*{\texorpdfstring{b) From your regression output,
calculate by how much the mean DBP changes for a
5kgm\textsuperscript{-2} increase in BMI? Can you verify this by
modifying your data and re-running your
regression?}{b) From your regression output, calculate by how much the mean DBP changes for a 5kgm-2 increase in BMI? Can you verify this by modifying your data and re-running your regression?}}\label{b-from-your-regression-output-calculate-by-how-much-the-mean-dbp-changes-for-a-5kgm-2-increase-in-bmi-can-you-verify-this-by-modifying-your-data-and-re-running-your-regression}
\addcontentsline{toc}{subsubsection}{b) From your regression output,
calculate by how much the mean DBP changes for a
5kgm\textsuperscript{-2} increase in BMI? Can you verify this by
modifying your data and re-running your regression?}

\textbf{Stata code and output}

\begin{Shaded}
\begin{Highlighting}[]
\CommentTok{/* Part a */}
\NormalTok{import delimited }\StringTok{"https://raw.githubusercontent.com/atpinto/RM1/refs/heads/main/Data/hers\_subset.csv"}\NormalTok{, }\KeywordTok{clear}
\KeywordTok{reg}\NormalTok{ dbp bmi}

\CommentTok{/* Part b*/}
\KeywordTok{gen}\NormalTok{ bmi5 = bmi / 5}
\KeywordTok{reg}\NormalTok{ dbp bmi5}
\NormalTok{\#\# (encoding automatically selected: ISO{-}8859{-}1)}
\NormalTok{\#\# (38 vars, 276 }\KeywordTok{obs}\NormalTok{)}
\NormalTok{\#\# }
\NormalTok{\#\# }
\NormalTok{\#\#       Source |       SS           df       MS      Number }\KeywordTok{of} \KeywordTok{obs}\NormalTok{   =       276}
\NormalTok{\#\# {-}{-}{-}{-}{-}{-}{-}{-}{-}{-}{-}{-}{-}+{-}{-}{-}{-}{-}{-}{-}{-}{-}{-}{-}{-}{-}{-}{-}{-}{-}{-}{-}{-}{-}{-}{-}{-}{-}{-}{-}{-}{-}{-}{-}{-}{-}{-}   }\FunctionTok{F}\NormalTok{(1, 274)       =      4.84}
\NormalTok{\#\#        Model |  423.883938         1  423.883938   Prob \textgreater{} }\FunctionTok{F}\NormalTok{        =    0.0286}
\NormalTok{\#\#     Residual |  23988.8842       274  87.5506722   R{-}squared       =    0.0174}
\NormalTok{\#\# {-}{-}{-}{-}{-}{-}{-}{-}{-}{-}{-}{-}{-}+{-}{-}{-}{-}{-}{-}{-}{-}{-}{-}{-}{-}{-}{-}{-}{-}{-}{-}{-}{-}{-}{-}{-}{-}{-}{-}{-}{-}{-}{-}{-}{-}{-}{-}   Adj R{-}squared   =    0.0138}
\NormalTok{\#\#        Total |  24412.7681       275  88.7737022   Root MSE        =    9.3569}
\NormalTok{\#\# }
\NormalTok{\#\# {-}{-}{-}{-}{-}{-}{-}{-}{-}{-}{-}{-}{-}{-}{-}{-}{-}{-}{-}{-}{-}{-}{-}{-}{-}{-}{-}{-}{-}{-}{-}{-}{-}{-}{-}{-}{-}{-}{-}{-}{-}{-}{-}{-}{-}{-}{-}{-}{-}{-}{-}{-}{-}{-}{-}{-}{-}{-}{-}{-}{-}{-}{-}{-}{-}{-}{-}{-}{-}{-}{-}{-}{-}{-}{-}{-}{-}{-}}
\NormalTok{\#\#          dbp | Coefficient  Std. err.      t    P\textgreater{}|t|     [95\% conf. interval]}
\NormalTok{\#\# {-}{-}{-}{-}{-}{-}{-}{-}{-}{-}{-}{-}{-}+{-}{-}{-}{-}{-}{-}{-}{-}{-}{-}{-}{-}{-}{-}{-}{-}{-}{-}{-}{-}{-}{-}{-}{-}{-}{-}{-}{-}{-}{-}{-}{-}{-}{-}{-}{-}{-}{-}{-}{-}{-}{-}{-}{-}{-}{-}{-}{-}{-}{-}{-}{-}{-}{-}{-}{-}{-}{-}{-}{-}{-}{-}{-}{-}}
\NormalTok{\#\#          bmi |   .2221827   .1009756     2.20   0.029     .0233961    .4209693}
\NormalTok{\#\#        }\DataTypeTok{\_cons}\NormalTok{ |   67.82592   2.923282    23.20   0.000     62.07097    73.58087}
\NormalTok{\#\# {-}{-}{-}{-}{-}{-}{-}{-}{-}{-}{-}{-}{-}{-}{-}{-}{-}{-}{-}{-}{-}{-}{-}{-}{-}{-}{-}{-}{-}{-}{-}{-}{-}{-}{-}{-}{-}{-}{-}{-}{-}{-}{-}{-}{-}{-}{-}{-}{-}{-}{-}{-}{-}{-}{-}{-}{-}{-}{-}{-}{-}{-}{-}{-}{-}{-}{-}{-}{-}{-}{-}{-}{-}{-}{-}{-}{-}{-}}
\NormalTok{\#\# }
\NormalTok{\#\# }
\NormalTok{\#\# }
\NormalTok{\#\#       Source |       SS           df       MS      Number }\KeywordTok{of} \KeywordTok{obs}\NormalTok{   =       276}
\NormalTok{\#\# {-}{-}{-}{-}{-}{-}{-}{-}{-}{-}{-}{-}{-}+{-}{-}{-}{-}{-}{-}{-}{-}{-}{-}{-}{-}{-}{-}{-}{-}{-}{-}{-}{-}{-}{-}{-}{-}{-}{-}{-}{-}{-}{-}{-}{-}{-}{-}   }\FunctionTok{F}\NormalTok{(1, 274)       =      4.84}
\NormalTok{\#\#        Model |  423.883889         1  423.883889   Prob \textgreater{} }\FunctionTok{F}\NormalTok{        =    0.0286}
\NormalTok{\#\#     Residual |  23988.8842       274  87.5506724   R{-}squared       =    0.0174}
\NormalTok{\#\# {-}{-}{-}{-}{-}{-}{-}{-}{-}{-}{-}{-}{-}+{-}{-}{-}{-}{-}{-}{-}{-}{-}{-}{-}{-}{-}{-}{-}{-}{-}{-}{-}{-}{-}{-}{-}{-}{-}{-}{-}{-}{-}{-}{-}{-}{-}{-}   Adj R{-}squared   =    0.0138}
\NormalTok{\#\#        Total |  24412.7681       275  88.7737022   Root MSE        =    9.3569}
\NormalTok{\#\# }
\NormalTok{\#\# {-}{-}{-}{-}{-}{-}{-}{-}{-}{-}{-}{-}{-}{-}{-}{-}{-}{-}{-}{-}{-}{-}{-}{-}{-}{-}{-}{-}{-}{-}{-}{-}{-}{-}{-}{-}{-}{-}{-}{-}{-}{-}{-}{-}{-}{-}{-}{-}{-}{-}{-}{-}{-}{-}{-}{-}{-}{-}{-}{-}{-}{-}{-}{-}{-}{-}{-}{-}{-}{-}{-}{-}{-}{-}{-}{-}{-}{-}}
\NormalTok{\#\#          dbp | Coefficient  Std. err.      t    P\textgreater{}|t|     [95\% conf. interval]}
\NormalTok{\#\# {-}{-}{-}{-}{-}{-}{-}{-}{-}{-}{-}{-}{-}+{-}{-}{-}{-}{-}{-}{-}{-}{-}{-}{-}{-}{-}{-}{-}{-}{-}{-}{-}{-}{-}{-}{-}{-}{-}{-}{-}{-}{-}{-}{-}{-}{-}{-}{-}{-}{-}{-}{-}{-}{-}{-}{-}{-}{-}{-}{-}{-}{-}{-}{-}{-}{-}{-}{-}{-}{-}{-}{-}{-}{-}{-}{-}{-}}
\NormalTok{\#\#         bmi5 |   1.110913   .5048781     2.20   0.029     .1169804    2.104847}
\NormalTok{\#\#        }\DataTypeTok{\_cons}\NormalTok{ |   67.82592   2.923282    23.20   0.000     62.07097    73.58087}
\NormalTok{\#\# {-}{-}{-}{-}{-}{-}{-}{-}{-}{-}{-}{-}{-}{-}{-}{-}{-}{-}{-}{-}{-}{-}{-}{-}{-}{-}{-}{-}{-}{-}{-}{-}{-}{-}{-}{-}{-}{-}{-}{-}{-}{-}{-}{-}{-}{-}{-}{-}{-}{-}{-}{-}{-}{-}{-}{-}{-}{-}{-}{-}{-}{-}{-}{-}{-}{-}{-}{-}{-}{-}{-}{-}{-}{-}{-}{-}{-}{-}}
\end{Highlighting}
\end{Shaded}

\textbf{R code and output}

\begin{Shaded}
\begin{Highlighting}[]
\CommentTok{\# Part a}
\NormalTok{hers\_subset }\OtherTok{\textless{}{-}} \FunctionTok{read.csv}\NormalTok{(}\StringTok{"https://raw.githubusercontent.com/atpinto/RM1/refs/heads/main/Data/hers\_subset.csv"}\NormalTok{)}
\NormalTok{lm.hers }\OtherTok{\textless{}{-}} \FunctionTok{lm}\NormalTok{(DBP }\SpecialCharTok{\textasciitilde{}}\NormalTok{ BMI, }\AttributeTok{data =}\NormalTok{ hers\_subset)}
\FunctionTok{summary}\NormalTok{(lm.hers)}
\DocumentationTok{\#\# }
\DocumentationTok{\#\# Call:}
\DocumentationTok{\#\# lm(formula = DBP \textasciitilde{} BMI, data = hers\_subset)}
\DocumentationTok{\#\# }
\DocumentationTok{\#\# Residuals:}
\DocumentationTok{\#\#      Min       1Q   Median       3Q      Max }
\DocumentationTok{\#\# {-}26.6420  {-}6.4584  {-}0.7538   5.8199  27.0639 }
\DocumentationTok{\#\# }
\DocumentationTok{\#\# Coefficients:}
\DocumentationTok{\#\#             Estimate Std. Error t value Pr(\textgreater{}|t|)    }
\DocumentationTok{\#\# (Intercept)  67.8259     2.9233    23.2   \textless{}2e{-}16 ***}
\DocumentationTok{\#\# BMI           0.2222     0.1010     2.2   0.0286 *  }
\DocumentationTok{\#\# {-}{-}{-}}
\DocumentationTok{\#\# Signif. codes:  0 \textquotesingle{}***\textquotesingle{} 0.001 \textquotesingle{}**\textquotesingle{} 0.01 \textquotesingle{}*\textquotesingle{} 0.05 \textquotesingle{}.\textquotesingle{} 0.1 \textquotesingle{} \textquotesingle{} 1}
\DocumentationTok{\#\# }
\DocumentationTok{\#\# Residual standard error: 9.357 on 274 degrees of freedom}
\DocumentationTok{\#\# Multiple R{-}squared:  0.01736,    Adjusted R{-}squared:  0.01378 }
\DocumentationTok{\#\# F{-}statistic: 4.842 on 1 and 274 DF,  p{-}value: 0.02862}
\FunctionTok{confint}\NormalTok{(lm.hers)}
\DocumentationTok{\#\#                   2.5 \%     97.5 \%}
\DocumentationTok{\#\# (Intercept) 62.07097446 73.5808680}
\DocumentationTok{\#\# BMI          0.02339609  0.4209693}

\CommentTok{\# Part b}
\NormalTok{hers\_subset}\SpecialCharTok{$}\NormalTok{BMI5 }\OtherTok{\textless{}{-}}\NormalTok{ hers\_subset}\SpecialCharTok{$}\NormalTok{BMI }\SpecialCharTok{/} \DecValTok{5}
\NormalTok{lm.hers }\OtherTok{\textless{}{-}} \FunctionTok{lm}\NormalTok{(DBP }\SpecialCharTok{\textasciitilde{}}\NormalTok{ BMI5, }\AttributeTok{data =}\NormalTok{ hers\_subset)}
\FunctionTok{summary}\NormalTok{(lm.hers)}
\DocumentationTok{\#\# }
\DocumentationTok{\#\# Call:}
\DocumentationTok{\#\# lm(formula = DBP \textasciitilde{} BMI5, data = hers\_subset)}
\DocumentationTok{\#\# }
\DocumentationTok{\#\# Residuals:}
\DocumentationTok{\#\#      Min       1Q   Median       3Q      Max }
\DocumentationTok{\#\# {-}26.6420  {-}6.4584  {-}0.7538   5.8199  27.0639 }
\DocumentationTok{\#\# }
\DocumentationTok{\#\# Coefficients:}
\DocumentationTok{\#\#             Estimate Std. Error t value Pr(\textgreater{}|t|)    }
\DocumentationTok{\#\# (Intercept)  67.8259     2.9233    23.2   \textless{}2e{-}16 ***}
\DocumentationTok{\#\# BMI5          1.1109     0.5049     2.2   0.0286 *  }
\DocumentationTok{\#\# {-}{-}{-}}
\DocumentationTok{\#\# Signif. codes:  0 \textquotesingle{}***\textquotesingle{} 0.001 \textquotesingle{}**\textquotesingle{} 0.01 \textquotesingle{}*\textquotesingle{} 0.05 \textquotesingle{}.\textquotesingle{} 0.1 \textquotesingle{} \textquotesingle{} 1}
\DocumentationTok{\#\# }
\DocumentationTok{\#\# Residual standard error: 9.357 on 274 degrees of freedom}
\DocumentationTok{\#\# Multiple R{-}squared:  0.01736,    Adjusted R{-}squared:  0.01378 }
\DocumentationTok{\#\# F{-}statistic: 4.842 on 1 and 274 DF,  p{-}value: 0.02862}
\end{Highlighting}
\end{Shaded}

We find evidence that diastolic blood pressure increases as body mass
index increases (P = 0.029). For every one kg/m\textsuperscript{-2}
increase in BMI, the mean diastolic blood pressure increases by
0.22mmHg, and we are 95\% confident the true increase lies between 0.023
and 0.42mmHg. BMI accounts for 1.7\% of the overall variability in
diastolic blood pressure.

If a 1 kg/m\textsuperscript{-2} increase in BMI accounts for a 0.22mmHg
increase in DBP, then a 5kg/m\textsuperscript{-2} increase in BMI
accounts for a \texttt{5x0.22\ =\ 1.1}mmHg increase in DBP. We can
confirm this in Stata or R by creating a new covariate \texttt{BMI5}
which is BMI scaled by a factor of 1/5 (so that a 1 increase in BMI5
corresponds to a 5 increase in BMI).

\subsubsection*{\texorpdfstring{c) Manually calculate the \(\beta_1\)
standard error, the t-value, p-value and
\(R^2\)}{c) Manually calculate the \textbackslash beta\_1 standard error, the t-value, p-value and R\^{}2}}\label{c-manually-calculate-the-beta_1-standard-error-the-t-value-p-value-and-r2}
\addcontentsline{toc}{subsubsection}{c) Manually calculate the
\(\beta_1\) standard error, the t-value, p-value and \(R^2\)}

From 3.3.7 of the textbook, the standard error of the regression
coefficient is as follows:

\(\hat{\text{se}}(\hat\beta_1) = \frac{\text{Root mean squared error}}{\hat\sigma_x \sqrt{(n-1)} }\)
We can use R or Stata to calculate \(\hat \sigma_x = 5.5879\).
Substituting this in we obtain
\(\hat{\text{se}}(\hat \beta_1) = 9.357/(5.5879 \sqrt{275}) = 0.101\) in
agreement with the Stata and R output.

The t-value is the regression coefficient divided by it's standard error
\(t = 0.222 / 0.101 = 2.2\), and the P-value can be calculated by
looking up a corresponding t-table with \(n-2 = 276-2 = 274\) degrees of
freedom:

Stata code and output for the p-value

\begin{Shaded}
\begin{Highlighting}[]
\CommentTok{/*Compute sigma\_x*/}
  \KeywordTok{summarize}\NormalTok{ bmi}
  \KeywordTok{display} \FunctionTok{r}\NormalTok{(}\FunctionTok{sd}\NormalTok{)}

\CommentTok{/*Compute pvalue*/}    
\NormalTok{  disp tprob(274,2.2)}
\NormalTok{\#\#     Variable |        Obs        Mean    Std. }\KeywordTok{dev}\NormalTok{.       Min        Max}
\NormalTok{\#\# {-}{-}{-}{-}{-}{-}{-}{-}{-}{-}{-}{-}{-}+{-}{-}{-}{-}{-}{-}{-}{-}{-}{-}{-}{-}{-}{-}{-}{-}{-}{-}{-}{-}{-}{-}{-}{-}{-}{-}{-}{-}{-}{-}{-}{-}{-}{-}{-}{-}{-}{-}{-}{-}{-}{-}{-}{-}{-}{-}{-}{-}{-}{-}{-}{-}{-}{-}{-}{-}{-}}
\NormalTok{\#\#          bmi |        276    28.40797    5.587877      18.05      46.49}
\NormalTok{\#\# }
\NormalTok{\#\# 5.5878775}
\NormalTok{\#\# }
\NormalTok{\#\# .0286418}
\end{Highlighting}
\end{Shaded}

R code and output

\begin{Shaded}
\begin{Highlighting}[]
\FunctionTok{sd}\NormalTok{(hers\_subset}\SpecialCharTok{$}\NormalTok{BMI) }\CommentTok{\#the sigma\_x}
\DocumentationTok{\#\# [1] 5.587877}
\NormalTok{(}\DecValTok{1}\SpecialCharTok{{-}}\FunctionTok{pt}\NormalTok{(}\FloatTok{2.2}\NormalTok{,}\DecValTok{274}\NormalTok{))}\SpecialCharTok{*}\DecValTok{2} \CommentTok{\#p{-}value}
\DocumentationTok{\#\# [1] 0.0286418}
\end{Highlighting}
\end{Shaded}

\(R^2\) is the fraction of the total variance explained by the model so
is equal to \(R^{2} = 423.88/24412.77 = 0.017\). These two variances are
default output in Stata. In R the model sum of squares and residual sum
of squares can be obtained using \texttt{anova(lm.hers)}, after which
the \(R^2\) can be calculated.

\subsubsection{\texorpdfstring{d) Based on your regression, make a
prediction for the mean diastolic blood pressure of people with a BMI of
28kgm\textsuperscript{-2}}{d) Based on your regression, make a prediction for the mean diastolic blood pressure of people with a BMI of 28kgm-2}}\label{d-based-on-your-regression-make-a-prediction-for-the-mean-diastolic-blood-pressure-of-people-with-a-bmi-of-28kgm-2}

\subsubsection{e) Calculate and interpret a confidence interval for this
prediction.}\label{e-calculate-and-interpret-a-confidence-interval-for-this-prediction.}

\subsubsection{f) Calculate and interpret a prediction interval for this
prediction.}\label{f-calculate-and-interpret-a-prediction-interval-for-this-prediction.}

Stata code and output

\begin{Shaded}
\begin{Highlighting}[]
\NormalTok{import delimited }\StringTok{"https://raw.githubusercontent.com/atpinto/RM1/refs/heads/main/Data/hers\_subset.csv"}
\KeywordTok{reg}\NormalTok{ dbp bmi}

\KeywordTok{lincom} \DataTypeTok{\_cons}\NormalTok{ + 28*bmi}

\KeywordTok{set} \KeywordTok{obs}\NormalTok{ 277}
\KeywordTok{replace}\NormalTok{ bmi = 28 }\KeywordTok{in}\NormalTok{ 277}
\KeywordTok{predict}\NormalTok{ fitDBP}
\KeywordTok{predict}\NormalTok{ seprDBP, stdf}
\KeywordTok{gen} \FunctionTok{upper}\NormalTok{ = fitDBP + 1.96*seprDBP }\KeywordTok{in}\NormalTok{ 277}
\KeywordTok{gen} \FunctionTok{lower}\NormalTok{ = fitDBP {-}1.96*seprDBP }\KeywordTok{in}\NormalTok{ 277}

\OtherTok{list}\NormalTok{ bmi fitDBP }\FunctionTok{lower} \FunctionTok{upper} \KeywordTok{in}\NormalTok{ 277}

\NormalTok{\#\# no; }\KeywordTok{data} \KeywordTok{in} \DecValTok{memory}\NormalTok{ would }\KeywordTok{be}\NormalTok{ lost}
\NormalTok{\#\# }\FunctionTok{r}\NormalTok{(4);}
\NormalTok{\#\# }
\NormalTok{\#\# }\FunctionTok{r}\NormalTok{(4);}
\end{Highlighting}
\end{Shaded}

R code and output

\begin{Shaded}
\begin{Highlighting}[]
\NormalTok{hers\_subset }\OtherTok{\textless{}{-}} \FunctionTok{read.csv}\NormalTok{(}\StringTok{"https://raw.githubusercontent.com/atpinto/RM1/refs/heads/main/Data/hers\_subset.csv"}\NormalTok{)}
\NormalTok{lm.hers }\OtherTok{\textless{}{-}} \FunctionTok{lm}\NormalTok{(DBP }\SpecialCharTok{\textasciitilde{}}\NormalTok{ BMI, }\AttributeTok{data =}\NormalTok{ hers\_subset)}

\NormalTok{new\_observation }\OtherTok{\textless{}{-}} \FunctionTok{data.frame}\NormalTok{(}\AttributeTok{BMI =} \DecValTok{28}\NormalTok{)}

\FunctionTok{predict}\NormalTok{(lm.hers, }
        \AttributeTok{newdata =}\NormalTok{ new\_observation, }
        \AttributeTok{interval=}\StringTok{"confidence"}\NormalTok{)}
\DocumentationTok{\#\#        fit      lwr      upr}
\DocumentationTok{\#\# 1 74.04704 72.93529 75.15878}

\FunctionTok{predict}\NormalTok{(lm.hers, }
        \AttributeTok{newdata =}\NormalTok{ new\_observation, }
        \AttributeTok{interval=}\StringTok{"prediction"}\NormalTok{)}
\DocumentationTok{\#\#        fit      lwr      upr}
\DocumentationTok{\#\# 1 74.04704 55.59306 92.50101}
\end{Highlighting}
\end{Shaded}

We predict that the mean diastolic blood pressure for those with a BMI
of 28kgm\textsuperscript{-2} to be 74mmHg. We are 95\% confident the
true mean lies between 72.9mmHg and 75.2mmHg. We expect that 95\% of
women with that BMI will have a diastolic blood pressure between
55.6mmHg and 92.5mmHg.




\end{document}
